% Part: history 
% Chapter: biographies 
% Section: julia-robinson
\documentclass[../../../include/open-logic-section]{subfiles}

\begin{document}

\olfileid[pt]{his}{bio}{rob}

\olsection{Julia Robinson}

\olphoto{robinson-julia}{Julia Robinson}

Julia Bowman Robinson foi uma matemática norte-americana. É conhecida
sobretudo por seu trabalho sobre problemas de decisão e, em especial,
por suas contribuições para a solução do décimo problema de Hilbert.
Robinson nasceu em St.~Louis, Missouri, em 8 de dezembro de 1919. Ela
se recordava de já se sentir intrigada pelos números quando criança
\citep[4]{Reid1986}. Aos nove anos, contraiu escarlatina e sofreu
várias crises recorrentes de febre reumática. Isso a obrigou a passar
grande parte do tempo na cama, causando-lhe atraso nos estudos. Embora
tenha conseguido recuperar esse atraso com a ajuda de professores
particulares, os efeitos físicos da doença tiveram um impacto
duradouro em sua vida.

Apesar das dificuldades da infância, Robinson concluiu o ensino médio
com vários prêmios em matemática e ciências. Iniciou seus estudos
universitários no San Diego State College e, no último ano,
transferiu-se para a Universidade da Califórnia em Berkeley. Ali, foi
influenciada pelo matemático Raphael Robinson. Os dois se tornaram
bons amigos e se casaram em 1941. Por ser esposa de um membro do corpo
docente, Robinson estava impedida de lecionar no departamento de
matemática de Berkeley. Embora continuasse a frequentar disciplinas
de matemática como ouvinte, esperava deixar a universidade e formar
uma família. Pouco depois do casamento, porém, Robinson contraiu
pneumonia. Disseram-lhe que havia um acúmulo considerável de tecido
cicatricial em seu coração devido à febre reumática que tivera na
infância. Dada a gravidade das cicatrizes, o médico previu que ela não
viveria além dos quarenta anos e a aconselhou a não ter filhos
\citep[13]{Reid1986}.

Robinson ficou deprimida por muito tempo, mas por fim decidiu
continuar a estudar matemática. Voltou a Berkeley e concluiu o
doutorado em 1948, sob a orientação de Alfred Tarski. Tarski havia
demonstrado que a teoria de primeira ordem dos números reais é
decidível, e do trabalho de G\"odel seguia que a teoria de primeira
ordem dos números naturais é indecidível. Era um importante problema
em aberto saber se a teoria de primeira ordem dos números racionais é
decidível. Em sua tese \citeyearpar{Robinson1949}, Robinson demonstrou
que não é.

Interessada em problemas de decisão, Robinson tentou então encontrar
uma solução para o décimo problema de Hilbert. Esse problema integrava
uma famosa lista de 23 problemas matemáticos proposta por David
Hilbert em 1900. O décimo problema pergunta se existe um algoritmo que
determine, em um tempo finito, se uma equação polinomial com
coeficientes inteiros, como $3x^2 - 2y +3 = 0$, tem uma solução nos
inteiros. Tais questões são conhecidas como \emph{problemas
diofantinos}. Após alguns sucessos iniciais, Robinson uniu esforços a
Martin Davis e Hilary Putnam, que também trabalhavam no problema. Eles
conseguiram mostrar que problemas diofantinos exponenciais (nos quais
as incógnitas também podem aparecer como expoentes) são indecidíveis
e mostraram que certa conjectura (mais tarde chamada de ``J.R.'')
implica que o décimo problema de Hilbert é indecidível
\citep{DavisPutnamRobinson1961}. Robinson continuou trabalhando no
problema durante toda a década de 1960. Em 1970, o jovem matemático
russo Yuri Matijasevich finalmente demonstrou a hipótese J.R. O
resultado combinado é hoje chamado de teorema de
Matijasevich--Robinson--Davis--Putnam, ou, abreviadamente, teorema
MRDP. Matijasevich e Robinson se tornaram amigos e colaboraram em
vários artigos. Em uma carta a Matijasevich, Robinson escreveu certa
vez que ``na verdade, fico muito contente que, trabalhando juntos (a
milhares de quilômetros de distância), estejamos obviamente fazendo
mais progresso do que qualquer um de nós conseguiria fazer sozinho''
\citep[45]{Matijasevich1992}.

Robinson foi a primeira mulher a presidir a Sociedade Americana de
Matemática e a primeira mulher eleita para a Academia Nacional de Ciências.
Morreu em 30 de julho de 1985, aos 65 anos, após receber o diagnóstico
de leucemia.

\begin{reading}
Os artigos matemáticos de Robinson estão disponíveis em suas
\textit{Obras reunidas} \citep{Robinson1996}, que também incluem uma
reimpressão de seu memorial biográfico da Academia Nacional de Ciências
\citep{Feferman1994}. Constance Reid, irmã mais velha de Robinson,
publicou uma ``Autobiografia de Julia'', baseada em entrevistas
\citep{Reid1986}, bem como uma biografia completa \citep{Reid1996}.
Um breve documentário sobre Robinson e o décimo problema de Hilbert
foi dirigido por George Csicsery \citep{Csicsery2016}. Para um breve
relato sobre as colaborações de Yuri Matijasevich com Robinson e a
influência dela sobre seu trabalho, veja \citep{Matijasevich1992}.
\end{reading}

\end{document}
